\honest{``PR'' is corrosive; ``reputation'' is not}{Anna Salamon}{2021}

This is ``a small detail,'' but worth writing up. If you try to ``navigate PR concerns,'' the idea will guide you ``toward harmful and confused actions.'' If you protect your ``reputation,'' ``brand'' or ``honor'' instead, ``I predict this will basically go fine.''

The difference, as I define it: honor means keeping fixed standards and being known to keep them, like the Lannisters in \textsc{Game of Thrones}, who always pay their debts. The gentleman's standard was not a single simple principle, but it was still a fixed, if unwritten, social standard that he was seen to keep. PR means having no fixed standards, and trying to predict what will upset ``people,'' especially the media or a self-reinforcing miasma, so as not to set them off. It is a ``weirder or loopier process,'' prone to self-reinforcing fears of shadows, reminiscent of ``Politics and the English Language'' and not of Strunk and White. \nb{Once honor is defined as having standards and PR as having none, the definitions do most of the work.}

My second example of honor is a gentleman of 1800 fighting a duel to defend it. \nb{Men died this way, and the post does not explain why protecting one's honor would ``basically go fine'' today.}

People who worry about PR think they need an outside expert, a ``PR consultant,'' whose advice they ``nervously heed.'' Honor is something a person ``can know or choose directly.'' So honor ``leaves them free to communicate something,'' while PR ``somehow (I think?) tends to pull a person away from communicating anything at all.''

For the main claim, that thinking about PR leads to harmful actions, I offer ``AFAICT,'' ``I predict'' and ``(I think?),'' and suggest you try it and see: whenever you find yourself navigating PR for yourself or a group, think instead of its honor, reputation or good name, and see whether everyone feels clearer, freer and more as though their feet are on the ground. I describe no person or organization that switched from thinking about PR to thinking about honor, or the other way.

My last line lists as related C. S. Lewis's ``The Inner Ring'' and ``The New York Times,'' a blog post by Rob Rhinehart. \nb{The post appeared on 14 February 2021. The day before, the New York Times had published a critical article about the blog Slate Star Codex and the rationalist community around it. Rhinehart, the founder of Soylent, says in his post that we are all ``controlled and oppressed and enslaved by the evil octopus of the New York Times.'' The post does not mention the Times article, or say whether it replies to it.}
